\pdfmapfile{+cm.map}
\pdfmapfile{+cmextra.map}
\documentclass[10pt,letterpaper]{article}
\usepackage[margin=.60in,headheight=14pt,headsep=13pt,footskip=23pt]{geometry}
\usepackage{fancyhdr,array,amsmath,amssymb,url,needspace}
\renewcommand{\familydefault}{\sfdefault}
\pagestyle{fancy}\fancyhf{}
\fancyhead[L]{\fontsize{10}{12}\selectfont Week 77 / Persistent holes / Adult guide}
\fancyfoot[L]{\fontsize{9}{10}\selectfont Bellingham Math Circle / Week 77 / F77-FAC-v4}
\fancyfoot[R]{\fontsize{9}{10}\selectfont\thepage}
\renewcommand{\headrulewidth}{0pt}\renewcommand{\footrulewidth}{0pt}
\setlength{\parindent}{0pt}\setlength{\parskip}{4pt}
\setlength{\tabcolsep}{5pt}\renewcommand{\arraystretch}{1.16}
\raggedright
\newcommand{\sectiontitle}[1]{\par\vspace{5pt}{\bfseries\fontsize{12}{14}\selectfont #1}\par\nopagebreak}
\newcommand{\itemtitle}[1]{\par\vspace{3pt}\textbf{#1}\ }
\newcommand{\hint}[1]{\par\textbf{Hints, in order.} #1}
\newcommand{\rim}[1]{\ensuremath{\mathrm{#1}}}
\begin{document}
\sectiontitle{Mathematical destination: facts and limits}
Work in a fixed finite planar triangulation. Triangle interiors do not overlap; pieces meet only along common edges or vertices. Add edges and filled faces without removing anything. Every face needs all three edges; a tied stage is recorded after its whole batch. All cancellation is modulo 2: two copies of one edge cancel.

\textbf{Boundary theorem.} A saved loop is gone precisely when some available filled triangles have its edge set as their combined boundary. Once gone, the same saved loop stays gone in every later growing board. On these planar boards a loop is gone exactly when every triangle inside it (inside either lobe of a figure eight) is filled, so survival needs no search over tile choices: a filling, if it exists, is unique, and the loop disappears when its last required triangle is filled. A loop's first appearance and final filling do \emph{not} by themselves identify a barcode interval for an arbitrary loop.

\textbf{Counts versus history.} The independent hole count is \(E-V+C-F\), where \(F\) counts filled triangles and \(C\) counts connected components. Counts alone do not determine survival. On the divided square write
\(P=AB+BC+AC,\ Q=AC+CD+DA,\ R=AB+BC+CD+DA\).
Then \(R=P+Q\). When \(R\) closes first (\(DA\) before \(AC\)), either one-tile filling leaves \(R\) nonzero; both kill it. When a triangle rim closes first, on the square or on two triangles meeting at a vertex, filling that triangle first kills the older loop. Problem 2 contains both histories. Problem 3 asks on which board every first loop can be gone at stage 5, with identical hole counts in every build.

\textbf{Different drawings, same class.} In Problem 6 four different edge loops represent the same nonzero class after \(ABO,BCO\) are filled. On a fresh unfilled board, their required tile sets differ. Exactly one pair permits either strict disappearance order. The proof uses which tiles are necessary, not guesses about which visible cavity inherits a name.

Children are meant to build, retain an unchanged edge record, test cancellation and find a construction or obstruction. A good first visit reaches one convincing square test and two contrasting histories. Adults support the completeness arguments in Problems 2 and 4 and the general reasoning in Problems 5 and 6. Formal homology, interval decomposition and general stability are adult extensions. A finite check is not an infinite proof; long-lived need not mean meaningful, or short-lived mean noise.

\sectiontitle{Who uses this prototype}
One shared Grades 4--5 packet, six student pages and six problems. Prerequisites: follow short edge labels, distinguish an edge from a filled triangle, keep a saved list unchanged, and cancel matching pairs with support. No multiplication or formal algebra is needed for the first route.

For the stated group, use one kit with the three older children and the organizer: two build/test while the third referees, then rotate each build. The K--1 and 2--3 tables either keep separately prepared activities or follow the adult-led route on the same pages (next page); decide before printing.

\sectiontitle{Print and prepare}
\textbf{Print, single-sided US Letter at 100\%.} For the 4--5 table: one shared student packet (six sheets), handed out one page at a time, and one spare packet; materials page 1 twice (a working and a spare set of tiles and strips) and page 2 once (\texttt{week-77-materials.pdf}, F77-M-v1, on card stock if available). For each pair at a younger table that uses this week: student page 2 (page 3 too at 2--3) and materials page 1. One copy of this guide for each adult.

\textbf{Materials page 1:} nine tiles at board scale, square \(ABC,ACD\); joined \(ABC,CDE\); fan \(ABO,BCO,CDO,DAO\); and the small \(XYZ\) launch tile matching student page 1. Each tile shows its name, its board and each edge's name along that side, so a child can read off the three tokens a tile adds. Seven labelled strips: the square's four 6.2\,cm sides and 8.8\,cm diagonal, and the joined board's \(AB,DE\) (5.8\,cm). Yarn cut to the strips' lengths also works. The fan starts with all eight printed edges; it needs no strips.

\textbf{Materials page 2 (one kit):} 33 edge tokens, three each of \(AB,BC,CD,DA,AC,CE,DE,AO,BO,CO,DO\); six shaded launch tokens, two each of \(XY,YZ,XZ\); nine blank tokens; stage cards \(0,2,3,4,5,8\) and two blanks. Three copies allow an edge from the saved loop and both incident tile boundaries to be on the table at once. Add two contrasting colored pencils, three envelopes, scissors and removable tape.

\textbf{Before children arrive:} cut along every line and keep the square, joined and fan tiles in separate envelopes: identically named triangles have different sizes. Check that each tile covers its printed triangle; fit-to-page printing shrinks every piece. On page 6 physically place \(ABO\) and \(BCO\) before starting: the printed board itself has no shaded fillings. Check that tokens can be moved without disturbing the saved list or board. Rehearse one whole cancellation and one reset with the real materials. \textbf{Physical rehearsal and classroom piloting have not been performed.}
\newpage
\sectiontitle{The hour}
\textbf{0--5 minutes:} running around. \textbf{5--8: shared launch}, for every table that uses this week (otherwise at the 4--5 table). Draw a large square with its diagonal on the board, or hold up a page 2 square with its strips and tiles. Lay the rim, add the diagonal, then cover one half and the other. Ask, ``Is every triangle inside the rim covered?'' Say, ``Gone means the tiles fill the saved rim. The edges stay.'' Do not turn the square into a list of answers to copy; children construct and test their own loops in Problem 1.

\textbf{4--5 table.} \textbf{8--10:} the token test on student page 1. Save \(XY,YZ,XZ\), place the \(XYZ\) tile after checking its three edges, add its three edge tokens, and cancel matching pairs. No tokens remain. Keep the saved list and board unchanged. Let the trio handle pieces and repeat the move. \textbf{10--20:} Problem 1.
\textbf{20--33:} Problem 2. \textbf{33--41:} Problem 3; on the joined board the answer comes quickly, so the time goes to the square's two cases.
\textbf{41--53:} choose Problem 4 or a short Problem 5 conversation; do not rush both.
\textbf{53--57:} share one saved loop and the evidence for its survival or death.
\textbf{57--60:} pack and record. Problem 6 merits a separate 15--25-minute continuation for a ready group.

\textbf{K--1 and 2--3 tables: adult-led, on the same pages, 8--55.} Each pair has student page 2's square, the square's tiles and five strips; no tokens. The adult reads and records; children place pieces. (1) Children lay strips to make a loop back to its starting dot; the adult writes its edges down. (2) Ask, ``Is every triangle inside your loop covered?'' Children add tiles one at a time, each only when its three edges are down, and say when the loop is filled. (3) Make another loop and compare. At 2--3, add Problem 2's question with page 3: ``Put down \(DA\) before \(AC\), or \(AC\) before \(DA\). Which loop closed first? Which tiles fill it?'' \emph{Key in plain words:} a loop around one triangle is filled as soon as that triangle is covered. The outside loop needs both tiles; one tile always leaves the other triangle inside it uncovered. The diagonal stays on the board even though it is not part of the outside loop. If \(DA\) closes the first loop, that loop is the outside rim and needs both tiles; if \(AC\) closes it first, it is \(ABC\)'s rim and tile \(ABC\) alone fills it. Whether younger children can keep a saved loop apart from the growing board is a prediction, not an observation.

\textbf{Stopping/fallback:} if labels or cancellation are still difficult, stay with Problem 1; skip full enumeration, stage changes and the fan. One child chooses a square loop, another predicts whether a single legal tile kills it, and the referee checks with tokens. Swap roles. Ask at sharing: \lq\lq What stayed the same in two builds, and what changed about the saved loop?''

\sectiontitle{Problem 1, page 2: one filling or two? \quad Core}
Children build from dots, save a loop, and test fillings on separate builds.
A one-tile example is \(P=AB,BC,AC\), killed by \(ABC\), or \(Q=AC,CD,DA\), killed by \(ACD\).
The loop surviving either single filling is \(R=AB,BC,CD,DA\).
Add \(AC\) before filling a half, even though it is not on the saved rim.

\textbf{Check with objects:} \(R+P=Q\) and \(R+Q=P\), so one tile cannot remove every token. Using both leaves none. These are the square's only three nonempty loops, up to starting point or direction. No new name for a visible cavity is needed.
\hint{1. Trace one tile's edge places. 2. Keep your saved list and try the other tile on a reset board. 3. Which edges remain after the doubled diagonal is removed?}

\sectiontitle{Problem 2, page 3: every schedule \quad Core}
Children choose both the edge order and the tile order. The complete four-case list is below; exactly three cases meet the target. Keep the first loop's original edge list throughout each replay. If enumeration stalls, find one successful build and move to Problem 3; leave completeness for a return visit.

\textbf{Complete answer and explanation.}
All rows start with \(AB,BC,CD\) at stage 0. A face entry means \lq\lq fill that triangle.''
\begin{center}
\begin{tabular}{ccccccl}
Stage 2 & Stage 4 & Stage 5 & Stage 8 & Saved loop & Gone at & Meets target\\
\hline
\(DA\)&\(AC\)&\(ABC\)&\(ACD\)&\(R\)&8&yes\\
\(DA\)&\(AC\)&\(ACD\)&\(ABC\)&\(R\)&8&yes\\
\(AC\)&\(DA\)&\(ABC\)&\(ACD\)&\(P\)&5&no\\
\(AC\)&\(DA\)&\(ACD\)&\(ABC\)&\(P\)&8&yes
\end{tabular}
\end{center}
Two choices for the first edge and two for the first tile give four cases, with everything else forced. If \(DA\) arrives first, both tiles are needed to fill its saved outer rim. If \(AC\) arrives first, the saved loop is exactly \(ABC\)'s boundary, so only that tile is needed. All four rows have hole counts \(0,1,2,1,0\) at stages \(0,2,4,5,8\).
\hint{1. Which edge did you add first? Save its completed loop before adding anything else. 2. Hold the edge order fixed and switch the two tiles. 3. After both tile choices, switch the edge order. Have you covered both choices at each decision?}

\needspace{12\baselineskip}
\sectiontitle{Problem 3, page 4: gone at stage 5? \quad Core}
Start with \(AC,BC,CD,CE\). The triangles are \(ABC\) and \(CDE\), meeting only at \(C\). Here is every schedule.
\begin{center}
\begin{tabular}{ccccc}
Stage 2 & Stage 4 & Stage 5 & Stage 8 & First loop gone at\\
\hline
\(AB\)&\(DE\)&\(ABC\)&\(CDE\)&5\\
\(AB\)&\(DE\)&\(CDE\)&\(ABC\)&8\\
\(DE\)&\(AB\)&\(ABC\)&\(CDE\)&8\\
\(DE\)&\(AB\)&\(CDE\)&\(ABC\)&5
\end{tabular}
\end{center}
If \(AB\) is first, save \(AB,BC,AC\); if \(DE\) is first, save \(CD,DE,CE\).
Each saved rim disappears exactly when its own tile arrives. The other tile has different edges, so cannot cancel it. \textbf{Yes, either way:} fill the first rim's own triangle at stage 5 (rows 1 and 4).

\textbf{On the square, only with \(AC\) first.} \(AC\) first saves \(ABC\)'s rim, gone at 5 when \(ABC\) is filled then (Problem 2, row 3). \(DA\) first saves the outer rim \(AB,BC,CD,DA\), which needs both tiles (Problem 1), so it lasts to stage 8 whichever tile comes first (Problem 2, rows 1--2).

All eight builds on the two boards have hole counts \(0,1,2,1,0\) at stages \(0,2,4,5,8\), yet the first loop is gone at 5 in some and at 8 in others. \textbf{No}, the counts alone do not tell when the first loop disappears.
\hint{1. Touch the triangle whose missing edge you closed first. Which tile fills that loop? 2. On the square, which loop do you save when \(DA\) comes first? Look back at Problem 1. 3. Find two builds with the same hole counts in which the first loop disappears at different stages.}

\sectiontitle{Problem 4, page 5: simultaneous arrivals \quad Further}
The printed schedule never has two holes at a finished stage. Its counts at \(0,2,4,8\) are \(0,1,1,0\): the edge and tile at 4 are a single recorded batch. There are four proposed one-item changes:
\begin{center}
\begin{tabular}{p{1.15in}p{.55in}p{4.7in}}
Move & To & Result\\
\hline
edge \(AC\)&3&legal; two holes after stage 3, before the tile at 4\\
edge \(AC\)&5&illegal; \(ABC\) at 4 would lack \(AC\)\\
tile \(ABC\)&3&illegal; its edge \(AC\) would still be absent\\
tile \(ABC\)&5&legal; two holes after stage 4, before the tile at 5
\end{tabular}
\end{center}
Thus exactly the first and fourth changes work. The rim still disappears at 8.
\hint{1. Freeze the board only after all of stage 4 is present. 2. Try moving the edge alone and check every tile's three sides. 3. Now move the tile alone. There are only two new times for each item.}

\sectiontitle{Problem 5, page 5: can a saved loop return? \quad Further}
\textbf{No.} Keep the particular selection of filled triangles that erased every token. All those triangles remain available at every later stage. Reusing exactly that selection still erases exactly the same saved edge list, even if extra edges and tiles have appeared. This proves the claim for every legal growing board; it does not depend on testing examples or on planarity.
\hint{1. Save the tiles you used when the loop disappeared. 2. Which of those tiles could be missing later? 3. Can you replay the identical token cancellation?}
A fresh reset is a new build; it does not contradict this statement.

\sectiontitle{Problem 6, page 6: two fixed loops \quad Return visit}
Put \(ABO,BCO\) tiles on the all-edge fan. A loop survives either next filling precisely when its required tile set includes \emph{both} \(CDO\) and \(DAO\). All four possible edge lists follow. Use any two distinct rows for the first request.
\begin{center}
\begin{tabular}{cll}
Label & Saved edges & Required tiles\\
\hline
\(L_0\)&\(AO,CO,CD,DA\)&\(CDO,DAO\)\\
\(L_1\)&\(AB,BO,CO,CD,DA\)&\(ABO,CDO,DAO\)\\
\(L_2\)&\(AO,BO,BC,CD,DA\)&\(BCO,CDO,DAO\)\\
\(L_3\)&\(AB,BC,CD,DA\)&\(ABO,BCO,CDO,DAO\)
\end{tabular}
\end{center}
In this starting board each row needs both remaining tiles, so it survives either first addition and disappears at the last one. The drawings are different; they represent the same current class because their differences are boundaries of already-filled \(ABO,BCO\).

\textbf{For the fresh-board request choose \(L_1,L_2\).}
Fill \(CDO,DAO,ABO,BCO\): \(L_1\) disappears at the third filling and \(L_2\) at the fourth.
Fill \(CDO,DAO,BCO,ABO\): \(L_2\) disappears at the third filling and \(L_1\) at the fourth. Their saved edge lists never change.

\textbf{Why this is the only reversible pair.} Each fan triangle has its own outside edge, appearing in no other tile. Hence the outside edges of a boundary force exactly which tiles were used; no alternative filling can erase a missing required tile. Including both unfilled tiles leaves only the four choices above for the two already-filled tiles. If one required set contains another, the smaller set is completed no later than the larger, in every order. Of the six pairs, only \(L_1,L_2\) have neither set containing the other.
\hint{1. Choose a loop and try each possible first filling on reset boards. 2. Which tiles are inside that loop? Can a missing one be replaced by tiles outside it? 3. On the fresh board, look for two loops that each require a tile the other does not.}
If children can give the two orders but cannot yet show completeness, keep that as a checked construction. The necessity argument is an adult-supported next step.

\sectiontitle{Adult mathematics: short proofs and next steps}
For a finite planar straight-line triangulation, the boundary map on face selections is injective. A nonempty finite union of selected triangles has an exposed edge on its outside boundary, which occurs once. Thus it cannot have zero mod-2 boundary. Two different fillings of the same loop would differ by a nonempty zero-boundary selection, a contradiction. The cycle space of a graph has dimension \(E-V+C\) (use a spanning forest); the \(F\) triangle boundaries are independent. Quotienting by them gives \(E-V+C-F\).

For the square with outer rim born at \(b\), diagonal at \(s\), and faces at \(u,v\), assume \(b<s\leq\min(u,v)\). The old class \(R=P+Q\) survives the first face and vanishes at the second. Choosing the first-filled triangle's rim as the additional basis vector gives intervals
\[
[b,\max(u,v))\quad\hbox{and}\quad[s,\min(u,v)).
\]
Omit a zero-length interval; at a death time the class is already gone. For Problem 4 the second interval is omitted in the original schedule, and becomes \([3,4)\) or \([4,5)\) in the two valid changes. On the joined board, filling older/newer first gives respectively \([2,5),[4,8)\) or \([2,8),[4,5)\). These conclusions use the actual unchanged loop classes across stages, not only hole counts. General persistence pairing is a later adult topic; not every saved loop is an interval generator.

\sectiontitle{Source lineage and what to record}
Edelsbrunner, Letscher and Zomorodian, \emph{Topological Persistence and Simplification} (2002), Section 2 (PDF pp.\ 2--3), Section 3 (pp.\ 4--6): mod-2 boundaries, growing complexes and persistence.
{\footnotesize\url{https://pub.ista.ac.at/~edels/Papers/2002-J-04-TopologicalPersistence.pdf}}
The square, joined-board and fan tasks and the elementary proofs here are independently devised teaching examples. Week 2's bonus challenges, Problem 1C (F02-BONUS-v1), find the same divided square's three do-nothing line sets, \(P\), \(Q\) and \(R\); this week adds filled triangles and time. No borrowed figures or source-paper exercises are included.

Afterward record which problems were reached, which moves children controlled, what required adult rescue, and one exact edge list, drawing or claim worth keeping. Mark each claim \emph{tried}, \emph{guessed}, \emph{checked in cases}, or \emph{explained}. Note material-fit and timing problems. Digital verification establishes the printed finite answers; no observed learning or classroom readiness is claimed.
\end{document}
